\documentclass[10pt,a4paper]{amsart}
\usepackage[margin=19mm]{geometry}
\usepackage{amsmath,amssymb,mathtools}
\usepackage[scr=rsfs,cal=euler]{mathalfa}
\usepackage{xfrac}
\usepackage[unicode,hidelinks,bookmarks=false]{hyperref}
\newcommand{\ie}{i.e.\ }
\newcommand{\cf}{cf.\ }
\newcommand{\resp}{respectively\ }
\newcommand{\Subsection}{\subsection}
\providecommand{\nobreakdash}{\nobreak\hbox{-}}
\setlength{\emergencystretch}{2em}
\multlinegap0pt
\usepackage{bm}
\usepackage{bbm}
\usepackage{dsfont}
\usepackage{xspace}
\usepackage{cancel}
\usepackage{mathtools}
\usepackage{amsthm}
\makeatletter
\@ifundefined{theorem}{\newtheorem{theorem}{Theorem}}{}
\@ifundefined{lemma}{\newtheorem{lemma}[theorem]{Lemma}}{}
\makeatother
\newcommand{\FF}{\mathbb{F}}
\newcommand{\MM}{\mathbb{M}}
\newcommand{\ZZ}{\mathbb{Z}}
\title[Dynamics of Dehn Twists in the Outer Automorphism Group of a]{Dynamics of Dehn Twists in the Outer Automorphism Group of a Free Group}
\author{Donggyun Seo}
\date{Source: arXiv:2511.15063v4 (CC BY 4.0)}
\begin{document}
\maketitle

\noindent\emph{Source and licence.} Dynamics of Dehn Twists in the Outer Automorphism Group of a Free Group. Version arXiv:2511.15063v4 (CC BY 4.0), distributed under the Creative Commons Attribution 4.0 International licence, as recorded in the arXiv licence metadata for this exact version and shown on the abstract page https://arxiv.org/abs/2511.15063v4. Full rights notice: licenses/arxiv-2511.15063-CC-BY-4.0.txt.\par\smallskip

\noindent\textit{Editorial scope.} Counted unit: the Lemma whose source label is \texttt{lem:circle\_types} in the article (source0.tex lines 724--728, source label \texttt{lem:circle\_types}), delivered together with the article's own complete original proof of it (source0.tex lines 730--826, one \texttt{proof} environment of 97 lines). No mathematics is written, repaired or re-derived: statement and proof are the article's own text line for line. Required context retained uncounted: none. Every definition, display and label of those lines is retained unchanged. Omitted from this package and not redistributed: 1--723, 729--729, 827--2531. Nothing inside a retained excerpt is omitted or altered: every retained body is delivered exactly as the article's lines are written, so no omission receipt of any kind accompanies this package. Citations: the retained excerpts carry no citation command at all, so no bibliography entry of the article is transcribed and no reference is redistributed. Reference adaptation: none was needed and none was made; every reference inside the delivered unit resolves inside it, so no unresolved reference or citation is produced. Printed slips kept literal: no printed word, symbol, hypothesis or number of the counted statement or of its proof was altered. Layout and class: the article's own class is \texttt{amsart} with packages including amsfonts, mathrsfs, amsthm, amsxtra, amssymb, amsmath, amscd, hyperref, inputenc, t1enc, eucal, indentfirst, graphicx, subfig; the wrapper is the lane's pinned \texttt{amsart} editorial wrapper at 10pt on A4, which restates the theorem-like environments the retained text needs and adds the article's own macro definitions that the retained text needs (\texttt{\textbackslash FF}, \texttt{\textbackslash MM}, \texttt{\textbackslash ZZ}), with extra wrapper packages (bm, bbm, dsfont, xspace, cancel, mathtools). The delivered numbering is the unit's own. Assets: no retained span contains a figure environment, a graphics-inclusion command, a PDF-page inclusion, a table or a TikZ picture; no image file is copied into this package, the package declares no asset dependency and no image file is redistributed with it. Licence: version arXiv:2511.15063v4 of this article is distributed under the Creative Commons Attribution 4.0 International licence, as recorded in the arXiv licence metadata for this exact version and shown on the abstract page https://arxiv.org/abs/2511.15063v4. A full rights notice is delivered at licenses/arxiv-2511.15063-CC-BY-4.0.txt. Characters: every retained excerpt is pure ASCII, so the delivered engine, XeLaTeX with its default Latin Modern text font, reports no missing character.
\begin{lemma} \label{lem:circle_types}
    Let $s$ and $t$ be essential surfaces of $\MM$, each a sphere or a torus.
    Then there are transverse representatives, again denoted $s$ and $t$, such that every circle of $s \cap t$ is null-homotopic in $\MM$.
    Consequently each circle of $s \cap t$ either bounds a disk in the surface containing it or, when that surface is a torus, is isotopic on it to the meridian.
\end{lemma}

\begin{proof}
    We make objects transverse wherever the corresponding intersections are intended to be transverse; in particular the final representatives $s$ and $t$ will be transverse to each other.

    If at least one of the two surfaces is a sphere, say $s$, then every circle of $s \cap t$ lies on $s$ and $\pi_1(s) = 1$, so it is null-homotopic in $\MM$ and there is nothing to prove.
    Assume therefore that $s$ and $t$ are essential tori.

    Let $u$ be an essential torus with core $\gamma_u$.
    Since $u$ is essential, the image of $\pi_1(u) \to \FF$ is nontrivial; since it is abelian and every abelian subgroup of a free group is cyclic, the image is infinite cyclic.
    Hence the kernel of $\pi_1(u) \to \FF$ is an infinite cyclic direct summand of $\pi_1(u) \cong \ZZ^2$; let $m \subset u$ be an embedded curve representing a generator of this kernel, that is, a meridian of $u$.

    As $\pi_1(u) \to \pi_1(\MM)$ is not injective and $u$ is two-sided, the loop theorem provides an embedded disk $D \subset \MM$ with $D \cap u = \partial D$ and with $\partial D$ an essential simple closed curve of $u$ lying in the kernel of $\pi_1(u) \to \FF$.
    An essential simple closed curve on a torus represents a primitive class of $\pi_1(u)$, so $\partial D$ represents $m^{\pm 1}$ and is isotopic to $m$ on $u$.
    Adjusting $D$ by this isotopy, we may assume
    \[
        D \cap u = \partial D = m .
    \]

    Compressing $u$ along $D$ yields an embedded sphere $\Sigma$.
    We claim that $\Sigma$ is essential.
    Suppose that $\Sigma$ bounds a ball $B$.
    In the standard compression picture there are two possibilities.
    If $D$ lies on the side of $\Sigma$ opposite $B$, then reversing the compression attaches a $1$--handle to $B$, producing a solid torus bounded by $u$, contrary to essentiality of $u$.
    If instead $D$ lies in $B$, then reversing the compression shows that $u \subset B$; hence every loop on $u$ is null-homotopic in $\MM$, again contrary to essentiality.
    Thus $\Sigma$ does not bound a ball.

    Reversing the compression exhibits $u$ as a sphere with a tube attached: there are disjoint open disks $\mathring{D}^{-}, \mathring{D}^{+} \subset \Sigma$ and an embedded arc $a$ with endpoints in $\partial D^{-} \cup \partial D^{+}$ and interior disjoint from $\Sigma$, such that
    \[
        u = \bigl(\Sigma \setminus (\mathring{D}^{-} \cup \mathring{D}^{+})\bigr) \cup A_u ,
        \qquad A_u \cong S^1 \times [0,1] ,
    \]
    where $A_u$ is the boundary annulus of a thin regular neighborhood of $a$, attached along $\partial D^{-}$ and $\partial D^{+}$.
    We call $A_u$ the \emph{tube} of $u$ and $a$ its \emph{spine}.
    The cross-section circles $S^1 \times \{r\}$ of $A_u$ are meridians of $u$, and $\gamma_u$ is represented by the loop obtained by closing $a$ through $\Sigma$.

    Write accordingly
    \[
        s = \bigl(\Sigma_s \setminus (\mathring{D}_s^{-} \cup \mathring{D}_s^{+})\bigr) \cup s_a ,
        \qquad
        t = \bigl(\Sigma_t \setminus (\mathring{D}_t^{-} \cup \mathring{D}_t^{+})\bigr) \cup t_a ,
    \]
    with spines $a_s, a_t$ and tubes $s_a, t_a \cong S^1 \times [0,1]$, and put $K_s := \Sigma_s \cup a_s$ and $K_t := \Sigma_t \cup a_t$.

    Since $\dim \MM = 3$ and the spines are $1$--dimensional, a small homotopy makes $a_s \cap a_t = \emptyset$.
    A further small homotopy makes $a_s$ transverse to $\Sigma_t$, $a_t$ transverse to $\Sigma_s$, and $\Sigma_s$ transverse to $\Sigma_t$, so that $a_s \cap \Sigma_t$ and $a_t \cap \Sigma_s$ are finite sets of interior points.
    The endpoints of $a_s$ may be slid freely on $\Sigma_s$, so we place them in $\Sigma_s \setminus K_t$, which is nonempty because $\Sigma_s \cap \Sigma_t$ is a union of circles on the sphere $\Sigma_s$ and $\Sigma_s \cap a_t$ is finite; symmetrically for the endpoints of $a_t$.

    Let $J_s \subseteq a_s$ be the closed subarc underlying the tube $s_a$, obtained by deleting small neighborhoods of the two endpoints of $a_s$; since those endpoints are disjoint from the compact set $K_t$, we may choose the deleted neighborhoods disjoint from $K_t$, so that
    \[
        a_s \cap K_t \subset \operatorname{int}(J_s) .
    \]
    Fix an identification $J_s \cong [0,1]$ compatible with the product structure $s_a \cong S^1 \times [0,1]$, and write
    \[
        \pi_s \colon s_a \cong S^1 \times [0,1] \longrightarrow [0,1] \cong J_s
    \]
    for the projection to the second factor, using the same symbol for the product projection of a regular neighborhood of $a_s$ onto $a_s$.

    For each $x \in a_s \cap K_t = a_s \cap \Sigma_t$ choose a closed interval $I_x \subset \operatorname{int}(J_s)$ with $x \in \operatorname{int}(I_x)$, the $I_x$ pairwise disjoint, and put
    \[
        I_s := \bigcup_x I_x \;\subsetneq\; J_s ,
    \]
    so that neither endpoint of $J_s$ lies in $I_s$.
    Then $a_s \setminus \operatorname{int}(I_s)$ is compact and disjoint from $K_t$, hence at positive distance from it for any fixed Riemannian metric on $\MM$.
    Choosing the regular neighborhoods of $K_s$ and $K_t$ inside sufficiently small metric neighborhoods, we obtain
    \[
        \pi_s\bigl(N(a_s) \cap N(K_t)\bigr) \subseteq I_s ,
    \]
    and since $t \subset N(K_t)$,
    \begin{equation} \label{eq:tube_projection_control}
        \pi_s(s_a \cap t) \subseteq I_s \subsetneq J_s .
    \end{equation}
    Shrinking further if necessary, we arrange the symmetric statement for $t_a \cap s$.

    After these isotopies, reconstruct the two sphere-with-tube surfaces using sufficiently thin regular neighborhoods, and continue to denote them by $s$ and $t$; they are isotopic to the original representatives through the preceding modifications.
    We choose the neighborhoods generically so that $s$ and $t$ are transverse.

    Let $c$ be a component of $s \cap t$ and suppose $[c] \neq 1$ in $\FF$.
    Since the image of $\pi_1(s) \to \FF$ is generated by $\gamma_s$, writing the slope of $c$ in a core--meridian basis of $s$ as $(p,q)$ gives $[c] = \gamma_s^{\,p}$; hence $p \neq 0$.

    After a small isotopy of the two boundary circles of $s_a$ we may assume they are transverse to $c$, so that $c \cap s_a$ is a disjoint union of properly embedded arcs and circles.
    Consider
    \[
        [\,c \cap s_a\,] \in H_1\bigl(s_a, \partial s_a\bigr) \cong \ZZ ,
    \]
    the isomorphism being given by algebraic intersection number with a meridian $S^1 \times \{r\}$.
    A circle in $s_a$, and an arc with both endpoints on the same boundary component of $s_a$, each have algebraic intersection number zero with a meridian, so only arcs joining the two boundary components contribute; the total is the algebraic intersection number of $c$ with a meridian of $s$, namely $p$.
    As $p \neq 0$, some component $\beta \subseteq c \cap s_a$ joins the two boundary components of $s_a$.
    Since the endpoints of $\beta$ lie on the two boundary components of $s_a$, its projection contains both endpoints of $J_s$; as $\pi_s(\beta)$ is connected, $\pi_s(\beta) = J_s$.
    On the other hand $\beta \subseteq c \subseteq t$, so $\beta \subseteq s_a \cap t$, and \eqref{eq:tube_projection_control} gives $\pi_s(\beta) \subseteq I_s \subsetneq J_s$.
    This contradiction shows $[c] = 1$.

    Let $c$ be a component of $s \cap t$ and let $u$ be either of the two surfaces, so that $c \subset u$.
    If $u$ is a sphere, then $c$ bounds a disk on $u$.
    Suppose $u$ is a torus.
    If $c$ is inessential on $u$, then $c$ bounds a disk on $u$.
    If $c$ is essential on $u$, then $c$ represents a primitive class of $\pi_1(u)$, of slope $(p,q)$ with $\gcd(p,q) = 1$ in a core--meridian basis, and its image in $\FF$ is $\gamma_u^{\,p}$.
    By the above this image is trivial, and $\gamma_u$ has infinite order, so $p = 0$; primitivity then forces $q = \pm 1$, so $c$ is isotopic on $u$ to the meridian.
\end{proof}

\end{document}
